\def\level{BTS}  % Classe à droite
\def\course{CIEL 2}
\def\chaptername{Lois continues} % Titre au centre
\def\docpart{Loi uniforme} % Partie du cours
\def\docname{C222}        % Identifiant à gauche

\pagestyle{cours_FP}
\makeheaderBTS{} % Affiche le bandeau

\begin{definition}[La loi uniforme]{}
Une \textbf{loi de probabilité} est dite \textbf{uniforme} si tous les éléments de l'univers ont la même probabilité d'occurrence.
	
\medskip

Une variable aléatoire $X$ suit une loi uniforme sur l'intervalle $\ff{a}{b}$ notée $X \sim \mathcal{U}\ff{a}{b}$ lorsque la densité de probabilité de $f$ est définir par:\\

\begin{minipage}{8cm}
$$f(x)=
\begin{cases}
\dfrac{1}{b-a} & \textrm{si} \qquad a \leqslant x \leqslant b \\
0 & \textrm{sinon}
\end{cases}$$

\end{minipage}
\begin{minipage}{8cm}
$$P(c \leqslant X \leqslant d)=\dfrac{d-c}{b-a}$$
\end{minipage}
\begin{center}
	\includegraphics[width = 0.6\linewidth]{C220_1.jpg}
	
\end{center}
\end{definition}


\begin{exercice_cours}[\quad]{}
Ci-contre est représentée la densité de probabilité de la loi uniforme X sur l'intervalle $\ff{1}{3}$:

\begin{minipage}{0.4\linewidth}

$$f(x)=
\begin{cases}
\dfrac{1}{2} & \textrm{si} \qquad 1 \leqslant x \leqslant 3 \\
0 & \textrm{sinon}
\end{cases}$$
\end{minipage}\hfill
\begin{minipage}{0.27\linewidth}
\includegraphics[width = 0.9\linewidth]{C220_2.jpg}
\end{minipage}\hfill
\begin{minipage}{0.33\linewidth}
	\begin{itemize}
		\item Calculer $P(X<2,5)$
		\item Calculer $P(X>2,5)$
		\item Calculer $P(1,5<X<2,5)$
		\item Calculer $P(1<X<2)$
	\end{itemize}
\end{minipage}
\end{exercice_cours}

\begin{propriete}[Propriétés des lois uniformes]{}
Soit $X$ une variable aléatoire suivant une loi uniforme sur l'intervalle $\ff{a}{b}$.

\begin{itemize}
	\item[$\bullet$] L'espérance de $X$ est : $E(X)=\dfrac{a+b}{2}$.
	\item[$\bullet$] La variance de $X$ est : $V(X)=\dfrac{(b-a)^2}{12}$.
	\item[$\bullet$] L'écart-type de $X$ est : $\sigma(X)=\dfrac{b-a}{\sqrt{12}}$.

\end{itemize}
\end{propriete}

\begin{exercice_cours}[A partir des données de l'exercice 1]{}
\begin{enumerate}
	\item Calculer l'espérance de $X$.\hfill\reponseLigne[7cm]{}.
	\item Calculer la variance de $X$.\hfill\reponseLigne[7cm]{}.
	\item Calculer l'écart-type de $X$.\hfill\reponseLigne[7cm]{}.
\end{enumerate}
\end{exercice_cours}

\pagebreak

\begin{exercice_cours}[\quad]{}

	$X$ est une variable aléatoire qui suit une loi uniforme sur l'intervalle $\ff{-2}{8}$.

\begin{enumerate}
	\item Donner la fonction de densité de $X$ 
	\begin{blocvisiblex}[height=1.5]
{\quad}
\end{blocvisiblex}
	\item Déterminer les probabilité suivantes:
	
	$$P(X  \in \ff{0}{3}) \qquad P(-1\leqslant X \leqslant 4) \qquad P(X\leqslant 5)  \qquad P(X=2)$$
	
\begin{blocvisiblex}[height=2.]
{\quad}
\end{blocvisiblex}	
	\item Déterminer l'espérance, la variance et l'écart-type de $X$.
	
\begin{blocvisiblex}[height=2.]
{\quad}
\end{blocvisiblex}	
	
\end{enumerate}
\end{exercice_cours}

\begin{exercice_cours}[\quad]{}
\bigskip

 Antoine arrive à l'arrêt du bus sans avoir consulté les horaires. Sur cette ligne, un bus part toutes les huit minutes. On note $X$ la variable aléatoire donnant, en minutes, le temps d'attente jusqu'au départ du bus.

\begin{enumerate}
	\item Quelle est la loi suivie par la variable aléatoire $X$?
	
\begin{blocvisiblex}[height=1.5]
{\quad}
\end{blocvisiblex}	
	\item Calculer la probabilité qu'Antoine attende:
	
	\begin{enumerate}
		\item moins de 2 minutes \hfill\reponseLigne[7cm]{}.
		\item entre 3 et 6 minutes \hfill\reponseLigne[7cm]{}.
		\item au plus 5 minutes \hfill\reponseLigne[7cm]{}.
		\item au moins 3 minutes \hfill\reponseLigne[7cm]{}.
		\item plus de 6 minutes \hfill\reponseLigne[7cm]{}.
	\end{enumerate}
	
	\item Calculer l'espérance de la variable aléatoire $X$ \hfill\reponseLigne[7cm]{}.
	\item Calculer la variance de la variable aléatoire $X$ \hfill\reponseLigne[7cm]{}.
	\item Calculer l'écart-type de la variable aléatoire $X$ \hfill\reponseLigne[7cm]{}.
\end{enumerate}
\end{exercice_cours}

\end{document}


\section{Lois discrètes}

\begin{propriete}[Lois discrètes]{}
Une \textbf{loi de probabilité} est dite \textbf{discrète} si elle est définie sur un ensemble fini ou dénombrable de valeurs.

On peut donc compter les éléments de l'univers et attribuer une probabilité à chacun d'eux. Si on jette un dé équilibré à six faces, la loi de probabilité associée est discrète car elle est définie sur l'ensemble fini $\{1, 2, 3, 4, 5, 6\}$.
\end{propriete}

\begin{exemple}[\quad]{}
	\begin{itemize}
		\item La \textbf{loi de Bernoulli} est une loi de probabilité discrète qui modélise une expérience aléatoire binaire, c'est-à-dire une expérience qui n'a que deux issues possibles: succès ou échec.
		\item La \textbf{loi binomiale} est une loi de probabilité discrète qui modélise le nombre de succès dans une séquence d'expériences de Bernoulli indépendantes.
	\end{itemize}
\end{exemple}


\section{Lois continues}

\begin{definition}[Lois continues]{}
Une \textbf{loi de probabilité} est dite \textbf{continue} si elle est définie sur un intervalle de valeurs réelles. 

\begin{itemize}
	\item les lois continues ne sont pas définies sur des points individuels,
	\item elles sont plutôt définies sur des intervalles de valeurs.
\end{itemize}
\end{definition}

\begin{exemple}[Exemples de lois continues]{}
	\begin{itemize}
		\item on définit une variable aléatoire qui associe à chaque train, le temps de trajet entre deux gares.
		\item on définit une variable aléatoire qui associe à chaque personne, sa taille.
		\item on définit une variable aléatoire qui associe à chaque jour, la température maximale enregistrée dans une ville.
	\end{itemize}
\end{exemple}

\begin{easybox}[Fonction de densité]{}
Une \textbf{fonction de densité} est une fonction qui permet de définir une loi de probabilité continue. Elle doit satisfaire les conditions suivantes:

\medskip

\begin{minipage}{0.7\linewidth}
	
\begin{itemize}
	\item $f(x) \geq 0$ pour tout $x \in \R$ et  $\displaystyle\int_{-\infty}^{+\infty} f(x) \, dx = 1$
	\item pour tout intervalle $\ff{a}{b}$, la probabilité que la variable aléatoire prenne une valeur dans cet intervalle est donnée par l'intégrale de la fonction de densité sur cet intervalle: \[P(X \in \ff{a}{b}) = P(a \leq X \leq b)= \int_{a}^{b} f(x) \, dx\]
\end{itemize}
\end{minipage}\hfill
\begin{minipage}{0.3\linewidth}
\includegraphics[width = 0.9\linewidth]{loiproba2.png}
\end{minipage}
\end{easybox}

\pagebreak

\begin{exemple}[Exemple de fonction de densité]{}
\begin{enumerate}
	\item Vérifier que la fonction $f$ définie ci-dessous est une fonction de densité.
	$$f(x) = \begin{cases}
2\e^{-x} & \text{si } x \in \ff{0}{\ln(2)}  \\
0 & \text{sinon}
\end{cases}$$
	\item Calculer $\displaystyle\int_{0}^{0,5} f(x) \, dx$ et interpréter le résultat en terme de probabilité.
	\item Calculer la probabilité que la variable aléatoire prenne une valeur supérieure à $0,5$.
	\item Calculer la probabilité que la variable aléatoire prenne une valeur comprise entre $0,2$ et $0,5$.
	\item Créer un curseur $a$ dans l'intervalle $\ff{0}{\ln(2)}$ et 
	\begin{itemize}
		\item déterminer $a$ tel que $P(X \leq a) = 0,25$,
		\item déterminer $a$ tel que $P(X \leq a) = 0,5$.
	\end{itemize}
\end{enumerate}
\end{exemple}

\begin{propriete}[Propriétés des lois continues]{}
\begin{itemize}
	\item $P(X=k)=\displaystyle\int_k^k\;f(x)\,dx=0$ 
	\item $P(X \leq a) = P(X < a)$ (la probabilité que la variable aléatoire prenne une valeur inférieure ou égale à $a$ est égale à la probabilité qu'elle prenne une valeur strictement inférieure à $a$)
	\item L'espérance d'une variable aléatoire continue $X$ est donnée par la formule: \[E(X) = \displaystyle\int_{-\infty}^{+\infty} x f(x) \, dx\]
\end{itemize}
\end{propriete}

\begin{exercice_cours}

Soit $X$ la variable aléatoire qui suit la loi de probabilité $f$ définie par $f(x)=0.015x-0.00075x^2$

\begin{enumerate}
	\item Montrer que $f$ est bien une densité de probabilité sur l'intervalle $\ff{0}{20}$
	
\begin{blocvisiblex}[height=2.]
{\quad}
\end{blocvisiblex}
	\item Calculer la probabilité de l'évènement \og $X$ est supérieure ou égale à 12 \fg{} 
	

\begin{blocvisiblex}[height=2.]
{\quad}
\end{blocvisiblex}	
	\item Calculer la probabilité de l'évènement \og $X$ est compris entre 5 et 15 \fg{} 
	

\begin{blocvisiblex}[height=2.]
{\quad}
\end{blocvisiblex}
\item Calculer l'espérance mathématique de $X$.


\begin{blocvisiblex}[height=2.]
{\quad}
\end{blocvisiblex}\end{enumerate}


\end{exercice_cours}